\section{Conclusion}


The design combines lattice-based threshold cryptography, TFHE-backed
homomorphic CRDT merge, non-custodial MPC wallets (CGGMP21 + FROST),
and three-mode regulatory disclosure. Every component is anchored in
Lux production code --- \texttt{luxfi/fhe}, \texttt{luxfi/lattice},
\texttt{luxfi/mpc}, \texttt{luxfi/age}, \texttt{luxfi/standard} --- with
formal verification artefacts in \texttt{luxfi/proofs}.

\bibliographystyle{plain}
\begin{thebibliography}{99}

\bibitem{shor1999}
P.~W.~Shor, ``Polynomial-time algorithms for prime factorization and
discrete logarithms on a quantum computer,'' \textit{SIAM Review},
vol.~41, no.~2, pp.~303--332, 1999.

\bibitem{nist-pqc}
NIST, ``Post-Quantum Cryptography Standardization,''
\url{https://csrc.nist.gov/projects/post-quantum-cryptography}, 2024.

\bibitem{nist-pqc-report}
NIST, ``Status Report on the Third Round of the NIST Post-Quantum
Cryptography Standardization Process,'' NISTIR~8413, 2022.

\bibitem{cggmp21}
R.~Canetti, R.~Gennaro, S.~Goldfeder, N.~Makriyannis, and U.~Peled,
``UC Non-Interactive, Proactive, Threshold ECDSA with Identifiable
Aborts,'' IACR ePrint 2021/060, 2021.

\bibitem{regev2005}
O.~Regev, ``On lattices, learning with errors, random linear codes,
and cryptography,'' \textit{STOC}, pp.~84--93, 2005.

\bibitem{gentry2009}
C.~Gentry, ``A fully homomorphic encryption scheme,'' Ph.D.\
dissertation, Stanford University, 2009.

\bibitem{lamport1979}
L.~Lamport, ``Constructing digital signatures from a one-way
function,'' Technical Report SRI-CSL-98, 1979.

\bibitem{rfc3526}
T.~Kivinen and M.~Kojo, ``More Modular Exponential (MODP)
Diffie-Hellman groups for Internet Key Exchange (IKE),'' RFC~3526,
2003.

\bibitem{ajtai1996}
M.~Ajtai, ``Generating hard instances of lattice problems (extended
abstract),'' \textit{STOC}, pp.~99--108, 1996.

\bibitem{micciancio-regev}
D.~Micciancio and O.~Regev, ``Lattice-based cryptography,'' in
\textit{Post-Quantum Cryptography}, Springer, pp.~147--191, 2009.

\bibitem{lpr2013}
V.~Lyubashevsky, C.~Peikert, and O.~Regev, ``On ideal lattices and
learning with errors over rings,'' \textit{Journal of the ACM},
vol.~60, no.~6, 2013.

\bibitem{gsw2013}
C.~Gentry, A.~Sahai, and B.~Waters, ``Homomorphic encryption from
learning with errors: Conceptually-simpler, asymptotically-faster,
attribute-based,'' \textit{CRYPTO}, 2013.

\bibitem{sphincs-plus}
D.~J.~Bernstein, A.~H\"{u}lsing, S.~K\"{o}lbl, R.~Niederhagen,
J.~Rijneveld, and P.~Schwabe, ``The SPHINCS+ Signature Framework,''
\textit{ACM CCS}, 2019.

\bibitem{tfhe2016}
I.~Chillotti, N.~Gama, M.~Georgieva, and M.~Izabach\`{e}ne, ``Faster
Fully Homomorphic Encryption: Bootstrapping in less than 0.1
Seconds,'' \textit{ASIACRYPT}, 2016.

\end{thebibliography}

